\begin{center}
\textbf{UNIVERSITY OF SCIENCE}\\
\textbf{ADVANCED PROGRAM IN COMPUTER SCIENCE}
\end{center}

\chapter*{Đề cương chi tiết}
\addcontentsline{toc}{section}{Thesis Proposal}

\begin{center}
\textbf{THESIS PROPOSAL}
\end{center}

\begin{tabular}{@{}L{3.6cm}L{11cm}@{}}
\textbf{Thesis title:} & RESEARCH AND BUILD A SYSSECOPS SYSTEM PROCESS FOR A HYBRID CLOUD ENVIRONMENT \\
\textbf{Thesis advisor:} & PhD. Tran Trung Dung, MSc. Chung Thuy Linh \\
\textbf{Students:} & Huynh Tuan Minh (22125055) -- Nguyen Chinh Thong (22125102) \\
\textbf{Type of thesis:} & Technology with demo application \\
\textbf{Duration:} & 30th September 2025 to 4th August 2026 \\
\end{tabular}

\bigskip
\noindent\textbf{Contents of thesis:}

\section*{1. Introduction}
\addcontentsline{toc}{subsection}{1. Introduction}

Traditional software development in hybrid cloud environments usually follows a linear process:
requirements, design, implementation, testing, and deployment. Although CI/CD automation has
accelerated testing and release, the coding and secure design stages remain costly and slow. Recent
advances in large language models (LLMs) make automatic code generation from natural language
practical, reducing development effort and delivery time.

However, AI-generated code is probabilistic and can contain logic errors, insecure patterns,
dependency risks, and policy violations. These risks become more critical in hybrid cloud
deployments, where workloads span public cloud and on-premises infrastructure, each with different
trust boundaries, network controls, and compliance requirements.

This thesis proposes a SysSecOps workflow that combines AI-assisted code generation with layered
security controls and automated operations. The workflow includes prompt governance, pre-execution
validation, static and dependency analysis, risk scoring, sandbox execution, runtime monitoring, and
automated incident response.

The architecture explicitly addresses:
\begin{itemize}
    \item Cross-environment identity and access control (cloud IAM + on-prem policy).
    \item Data and artifact movement between cloud and on-prem zones.
    \item Consistent security checks across heterogeneous runtime environments.
    \item Failover and observability across distributed infrastructure.
\end{itemize}

\section*{2. Objectives}
\addcontentsline{toc}{subsection}{2. Objectives}

\textbf{Primary objective:}

The primary objective of this research is to design and develop a SysSecOps operational process for
hybrid cloud environments that integrates system administration, continuous security, and software
delivery into a unified workflow, improving the coordination between infrastructure management,
security operations, and software development while supporting secure, automated, and efficient
deployment across heterogeneous cloud infrastructures.

\textbf{Specific objectives:}

\begin{enumerate}
    \item Investigate the operational and security challenges of hybrid cloud environments and
    analyze the limitations of existing approaches to integrating system administration, software
    development, and continuous security.
    \item Design a SysSecOps operational process that integrates infrastructure management, automated
    security assessment, policy enforcement, and deployment into a single, continuous-security
    workflow.
    \item Develop a prototype SysSecOps system implementing the proposed process, integrating
    multiple security assessment tools, infrastructure automation, and continuous security mechanisms
    within a hybrid cloud environment.
    \item Establish a unified security assessment mechanism that aggregates findings from
    heterogeneous security tools -- including a learned code-risk signal -- and supports risk
    evaluation to prioritize remediation.
    \item Evaluate the prototype through representative hybrid-cloud deployment scenarios, examining
    operational feasibility, applicability to an existing on-premises target, and the correctness of
    the unified risk-evaluation mechanism.
\end{enumerate}

\section*{3. Scope of the Thesis}
\addcontentsline{toc}{subsection}{3. Scope of the Thesis}

\textbf{Included scope:}
\begin{itemize}
    \item An end-to-end prototype covering AI-assisted generation, synthesis/validation, the
    multi-scanner security gate, risk scoring, deployment governance, and post-deployment monitoring.
    \item A hybrid testbed consisting of an AWS account (public side) and on-premises Linux nodes
    (private side) joined by a secure mesh-VPN connection.
    \item Security validation of AI-generated application code and IaC artifacts.
    \item A learned code-risk component (logistic regression over Bandit/Semgrep features) as one
    input to the risk-scoring engine.
\end{itemize}

\textbf{Excluded scope:}
\begin{itemize}
    \item Training large language models from scratch.
    \item Enterprise-scale, multi-account, or multi-tenant production deployment.
    \item Formal, statistically powered benchmarking against a manual-operations baseline.
\end{itemize}

\section*{4. System Architecture and Technology Stack}
\addcontentsline{toc}{subsection}{4. System Architecture and Technology Stack}

The system is organised into three cooperating zones. Zone~1 produces IaC; Zone~2 evaluates and
gates it; Zone~3 provisions the hybrid infrastructure and continuously observes it. The
risk-scoring engine sits at the boundary of Zone~1 and Zone~2, consuming scanner output and emitting
the decision that governs whether code proceeds to deployment or returns to Zone~1 for
regeneration.

\bigskip
\noindent\textit{The architecture diagram is the same three-zone diagram presented in Chapter~3
(Figure~\ref{fig:arch-overview}). It shows the three zones and the principal flow of control between
them: the solid arrows denote the forward path of a change from prompt to deployed and monitored
infrastructure, and the dashed arrow denotes the remediation feedback that returns rejected
generated code to Zone~1.}

\bigskip
\noindent\textbf{AI and orchestration:}
\begin{itemize}
    \item OpenAI API (LLM-based code generation)
    \item Python orchestration service (FastAPI)
\end{itemize}

\noindent\textbf{Security pipeline:}
\begin{itemize}
    \item Bandit, Semgrep, Checkov, cfn-lint, ansible-lint, infracost
    \item OWASP Dependency-Check (optional secondary validation)
\end{itemize}

\noindent\textbf{Execution and deployment:}
\begin{itemize}
    \item On-prem environment: VMware/VirtualBox-based Linux nodes
    \item Terraform/CDK for reproducible infrastructure provisioning
\end{itemize}

\noindent\textbf{Monitoring and response:}
\begin{itemize}
    \item Prometheus (on-prem workloads)
    \item Alertmanager/webhook automation for quarantine, rollback, or block actions
\end{itemize}

\section*{5. Methods}
\addcontentsline{toc}{subsection}{5. Methods}

This research adopts a design-and-evaluation methodology to develop and validate the SysSecOps
workflow.

\textbf{1. System Design}

A modular architecture will be designed to integrate AI-based code generation with multiple
security layers, including validation, analysis, and monitoring components.

\textbf{2. Implementation}

A prototype system will be developed using appropriate tools and technologies, such as large
language model APIs for code generation, static analysis tools for vulnerability detection.

\textbf{3. Security Analysis Techniques}

The system will incorporate:
\begin{itemize}
    \item Static code analysis for identifying known vulnerability patterns.
    \item Dependency and library scanning for supply chain risks.
    \item Risk scoring algorithms to prioritize potential threats.
\end{itemize}

\textbf{4. Experimental Setup}

The workflow is exercised through representative pass, review, reject, and degraded-component
(missing scanner/credential) scenarios on the public-cloud (CDK) side, and through configuration of
an existing on-premises virtual machine over a mesh-VPN connection on the private (Ansible) side.

\textbf{5. Evaluation Metrics}

The system will be evaluated based on:
\begin{itemize}
    \item Operational feasibility and reliability: completion rate, stage-level completion,
    unexpected failures versus valid security rejections, and graceful degradation.
    \item Hybrid infrastructure applicability: connectivity, assessment, gate enforcement,
    configuration application, and state verification for an existing on-premises node.
    \item Unified risk-evaluation capability: accuracy, precision, recall, F1-score, and ROC--AUC of
    the learned code-risk component, plus correctness of score aggregation and decision thresholds.
\end{itemize}

\textbf{6. Limitations}

Despite the effectiveness of the proposed SysSecOps workflow in mitigating risks associated with
AI-generated code, several limitations remain:

\noindent\textit{Dependence on LLM Training Data:}
\begin{itemize}
    \item The model may produce outdated practices, insecure patterns, or suboptimal designs if such
    patterns exist in its training corpus.
    \item The model may not fully capture domain-specific requirements or contextual constraints.
\end{itemize}

\noindent\textit{False Positives vs. False Negatives Trade-off:}
\begin{itemize}
    \item Static analysis and security scanning tools (e.g., Checkov, Semgrep) may generate false
    positives, identifying issues that are not actual vulnerabilities.
    \item Reducing false positives by relaxing detection rules can increase the risk of false
    negatives, where real vulnerabilities are missed.
\end{itemize}

\noindent\textit{System Overhead:}
\begin{itemize}
    \item Additional computational overhead, increasing build time and resource consumption compared
    to traditional pipelines.
\end{itemize}

\noindent\textit{Incomplete Coverage:}
\begin{itemize}
    \item Not guarantee that all execution paths are fully tested or validated, even with automated
    testing mechanisms.
\end{itemize}

\section*{6. Expected Results}
\addcontentsline{toc}{subsection}{6. Expected Results}

The proposed workflow is expected to:
\begin{itemize}
    \item Reduce secure delivery time compared with traditional development pipelines.
    \item Increase vulnerability detection coverage for AI-generated code and IaC.
    \item Reduce production risk by runtime response automation.
    \item Demonstrate stable and policy-consistent operation across hybrid cloud infrastructure.
\end{itemize}

\textbf{Final deliverables:}
\begin{enumerate}
    \item A working SysSecOps prototype for AI-generated code.
    \item A hybrid cloud experimental environment definition (CDK app, Ansible playbook).
    \item Evaluation dataset and benchmark scenarios.
    \item Comparative performance and security analysis report.
\end{enumerate}

\section*{Research Timelines}
\addcontentsline{toc}{subsection}{Research Timelines}

\begin{table}[htb!]
\centering
\caption{Research timeline.}
\label{tab:proposal-timeline}
\begin{tabular}{|L{5.6cm}|C{2.2cm}|C{2.0cm}|C{2.4cm}|}
\hline
\textbf{Tasks} & \textbf{Time} & \textbf{Duration} & \textbf{Individual} \\
\hline
Study AWS cloud infrastructure & 9/2025 & 3 months & Minh, Thong \\
\hline
Design pipeline & 1/2026 & 1 month & Minh \\
\hline
Select AI model & 1/2026 & 1 month & Thong \\
\hline
Select service & 3/2026 & 1 week & Thong \\
\hline
Configure service, IAM policies, secure API & 15/3/2026 & 1 week & Thong \\
\hline
Draw and fix pipeline & 15/3/2026 & 1 week & Minh, Thong \\
\hline
Implement pre-validation & 15/3/2026 & 2 weeks & Minh \\
\hline
Bring pre-validation to service & 29/3/2026 & 1 week & Minh, Thong \\
\hline
Implement security analysis & 5/4/2026 & 2 weeks & Minh \\
\hline
Setup sandbox execution & 5/4/2026 & 2 weeks & Thong \\
\hline
Bring security analysis to service & 19/4/2026 & 1 week & Minh, Thong \\
\hline
Implement risk scoring & 5/2026 & 1 month & Minh \\
\hline
Setup runtime monitoring & 5/2026 & 1 month & Thong \\
\hline
Write paper & 6/2026 & 2 months & Minh, Thong \\
\hline
Evaluate and finalize paper & 8/2026 & 1 week & Minh, Thong \\
\hline
\end{tabular}
\end{table}

\section*{Reference}
\addcontentsline{toc}{subsection}{Reference}

\begin{enumerate}
    \item R. Santos, A. Correia, and J. P. Sousa, ``Security Best Practices for Cloud-Native
    Applications: A Systematic Review,'' \textit{IEEE Access}, vol.~11, pp.~45230--45245, 2023.
    \item HashiCorp, ``Terraform: Infrastructure as Code,'' HashiCorp Documentation, 2024. [Online].
    Available: \url{https://developer.hashicorp.com/terraform/docs}.
    \item OWASP Foundation, ``OWASP Top 10,'' 2021 (vulnerability taxonomy baseline).
    \item Shetti, Milan, \textit{Nubes: Towards building a Secure and Scalable Hybrid Cloud
    Infrastructure} (University of Minnesota thesis, 2021).
    \item Shaileshbhai, Khushi, \textit{Security Challenges and Mitigation Strategies in Hybrid Cloud
    Environments} (January 31, 2025). Available at SSRN: \url{https://ssrn.com/abstract=5118715} or
    \url{http://dx.doi.org/10.2139/ssrn.5118715}.
\end{enumerate}

\bigskip
\begin{tabular}{@{}L{7.5cm}L{7.5cm}@{}}
Approved by the advisor & Ho Chi Minh city, 06/04/2026 \\[24pt]
\hrulefill & \hrulefill \\
Signature of advisor & Signature(s) of student(s) \\
\end{tabular}